\chapter{生きたニューロンでできたマシン}
\chaptermark{生きたニューロンでできたマシン}
\label{sec:livingmachine}
\begin{chaptergoals}
\item なぜブロードキャストチャネルのない培養は一斉発火を起こし、その間隔は何で決まるのか；
\item 生きたネットワークがフィードバックで応答を変える仕組みと、まだ証明されていないこと；
\item オルガノイドが、線形読み出しを外側のシリコンで学習させるリザバーになる仕組み；
\item 培養がこれまでどう教えられてきたか、なぜ教師がまだシャーレの外にいるのか。
\end{chaptergoals}

\section{回路図の見えるテストベンチ}

第\ref{sec:debugging}章では、回路図のないマシンをデバッグした。プローブが拾えるのは800億個のニューロンのうち1000個であり、残りはすべて推測するしかない。この立場に置かれたエンジニアなら、単純なことをするだろう。回路の小さな一片をブレッドボードに組み、部品を一つ残らず見えるようにするのである。ニューロンでも同じことができる。

大脳皮質のニューロンをばらばらにし、電極アレイを載せたチップの上で育てる。電極の数は数十から数万である。数週間のうちに細胞は突起を伸ばし、数千から数十万個のニューロンからなるネットワークにつながる。大部分は\nt{GLUT}で、少数の\nt{GABA}が混じる。その割合は標本によって異なり、海馬の培養ではニューロンの約 $6\%$ である\cite{Benson1994}。各電極は、第\ref{sec:debugging}章のプローブのように聴くこともでき、テスト信号のように電流を流すこともできる。テストベンチのもう一つの種類が\emph{オルガノイド}である。ヒトの幹細胞から育てた、大きさ1ミリほどの三次元の神経組織の塊である\cite{Lancaster2013}。こうしたテストベンチの上で、生きたネットワークは何か月も動き続ける。オルガノイドの実験をネットワーク越しに遠隔で100日以上続けて行える施設もある\cite{Jordan2024}。この分野は\emph{生物学的計算}、あるいは\emph{オルガノイド知能}と呼ばれる\cite{Smirnova2023}。

テストベンチには脳との大きな違いが二つある。一つは小ささで、ニューロンの数は脳の10万分の1である。もう一つは、ブロードキャストチャネルがないことである。大脳皮質から作った培養には\nt{DOPA}、\nt{SERO}、\nt{NORA}、\nt{ACET}のニューロンがない。データバスとフロー制御はあるが制御レジスタのないマシンである。これに何ができるかを見ていこう。

\section{ネットワークが自分でしていること}

放っておかれた培養は、黙ってもいないし、一様にノイズを出してもいない。ときどきネットワークのほぼ全体が一度に発火し、つまり\emph{一斉発火}し、また静まる。一斉発火の間隔は、培養によって1秒から数分である\cite{Wagenaar2006}。エンジニアには見慣れた光景である。強い正帰還をもち負荷のない増幅器は発振器になる。

仕組みはすでに知っている。\nt{GLUT}ニューロン同士の強い相互結合は、第\ref{sec:gaba}章で命題\ref{prop:ei}の条件が破れたときのように雪崩を起こす。雪崩はシナプスの伝達物質の蓄えを急速に使い果たし\eqref{eq:depression}、ネットワークは静かになる。蓄えは時定数 $\tau_{\text{rec}}$ で回復し、新しい雪崩に足りる割合 $x^*$ まで戻ると、次の一斉発火が起こる。一斉発火の間隔は
\begin{empheq}[box=\keybox]{equation}
T_{\text{burst}} \;\approx\; \tau_{\text{rec}}\,\ln\frac{1}{1-x^*} .
\label{eq:burstinterval}
\end{empheq}
第\ref{sec:code}章の $\tau_{\text{rec}}=0.8\,\text{s}$ と $x^*=0.9$ ではおよそ2秒、$x^*=0.99$ ではおよそ4秒になる。これは本書の $\tau_{\text{rec}}$ を使ったモデルの推定であり、測定された間隔の短い側に入る。微小培養での直接測定では、一斉発火のあと伝達物質の蓄えが回復するのに数十秒かかる\cite{CohenSegal2011}。つまりそこでは $\tau_{\text{rec}}$ がもっと大きい。その $\tau_{\text{rec}}$ を使えば、式\eqref{eq:burstinterval}は数十秒から数分を与え、遅い培養と合う。これは弛張型の発振器であり、第\ref{sec:muscles}章のリズム発生器の親戚である。あちらでは群の疲労が発振器を再充電したが、ここでは伝達物質の蓄えがその役をする。

一斉発火は、ネットワークがすることのないときにすることである。生きた脳では、似た光景が深い睡眠（第\ref{sec:sleep}章）と、本物の入力が入ってくる前の発達中の脳に見られる。この類似は示唆的だが、仕組みが同じだというのは仮説である。

\section{ドーパミンの教師なしでネットワークを教える}

培養には\nt{DOPA}がないので、「良いか悪いか」の信号は自分たちで、電極から与えるしかない。実験によれば、テストベンチ上のネットワークは確かに応答を変える。そこで最も興味深いのは、\emph{何で}教えるかである。

\emph{黙る教師。}1本の電極から1〜3秒に1回ネットワークを刺激し、別の電極に刺激後 $50\pm10\ms$ で目的の応答が現れるまで続ける。応答が現れたら、刺激を5分間止める。このサイクルを数回繰り返すと、刺激に対してすぐに目的の応答が出るようになる\cite{Shahaf2001}。ここでの報酬は静けさである。刺激はネットワークを揺さぶり、刺激の停止はネットワークを正しい応答を出した状態のまま残す。これは第\ref{sec:dofa}章のランダムな摂動による勾配降下（命題\ref{prop:gradient}）であり、揺さぶられていること自体が誤差の役をする。正しくなるまで揺さぶるのである。

\emph{テニスをするネットワーク。}電極が密に並んだテストベンチの実験で、約100万個のニューロンが、ラケット1本のコンピュータテニスゲームにつながれた\cite{Kagan2022}。ボールの位置は「感覚」電極に電流として与えられた。ボールがどこにあるかは場所で、どれだけ遠いかは頻度で符号化された。二つの「運動」領域の活動がラケットを上下に動かした。学習の規則は変わっていた。ラケットがボールを打ち返すと、ネットワークは短い\emph{予測可能な}刺激を受けた。ミスすると、数秒間の\emph{予測不能な}ランダム電流を受けた。20分のプレイでラリーは長くなった。セッション内の有意な伸びは完全なフィードバックのときだけだった。刺激の代わりに無音にした場合も、セッションの終わりにはフィードバックなしより上手にプレイしたが、効果は小さかった。そして数日にわたるセッション間の安定した学習はなかった。この実験の詳しい分析は第\ref{sec:dishbrain}章にある。

著者らは、ネットワークは自分の入力が予測可能になるように行動する、という仮説でこの結果を説明した\cite{Friston2010}。第\ref{sec:code}章のスパイクの節約や第\ref{sec:recognition}章のフレームの予測と同じ考えである。入力が予想どおりなら、マシンの消費は少なくて済む。しかし実験そのものも、とりわけ培養の「知性」についての著者らの言葉も、厳しい批判を招いた。効果は小さく、もっと単純にも説明できるというのである\cite{Balci2023}。公正な評価はこうである。テストベンチ上の培養がフィードバックによって振る舞いを変えることは測定済みである。それがまさに自分の入力を予測するように学習しているというのは仮説である。

\emph{体を動かすネットワーク。}同じように、培養を単純な仮想の「体」につなぎ、訓練刺激の時空間パターンを選ぶことで、体を目標へ動かすよう教えた実験もある\cite{Bakkum2008}。これらの実験のどれにもドーパミンはない。教師の役をするのは、エンジニアが選ぶ電流のパターンである。教師がプログラムやドーパミンそのものになると何が変わるかは、\ref{sec:dishdopamine}--\ref{sec:cartpole}節で見る。

\begin{figure}[t]
\centering
\begin{tikzpicture}[>=Stealth,
  blk/.style={draw,very thick,rounded corners=2pt,align=center,font=\footnotesize,minimum height=0.8cm,fill=blue!5}]
% (a) closed loop
\node[font=\small] at (1.5,4.3) {(a) 閉ループ};
\node[blk,text width=2.6cm] (G) at (1.5,3.3) {ゲーム：\\ボールとラケット};
\draw[very thick,fill=orange!10,rounded corners=3pt] (0.5,0.2) rectangle (2.5,1.6);
\foreach \x in {0.7,1.0,...,2.35}{\foreach \y in {0.4,0.7,1.0,1.3}{\fill[gray!60] (\x,\y) circle (0.04);}}
\draw[->,thick,blue!60!black] (G.west) -| (-0.5,0.9) -- (0.5,0.9);
\node[font=\scriptsize,blue!60!black,anchor=east] at (-0.6,2.1) {ボール $\to$ 電流};
\draw[->,thick,red!70!black] (2.5,0.9) -- (3.5,0.9) |- (G.east);
\node[font=\scriptsize,red!70!black,anchor=west] at (3.6,2.1) {スパイク $\to$ ラケット};
\node[font=\scriptsize] at (1.5,-0.2) {電極アレイ上のニューロン};
\node[font=\scriptsize,align=center] at (1.5,-0.85) {ヒット：予測可能な刺激\\ミス：ランダム電流};
% (b) reservoir
\begin{scope}[xshift=7.9cm]
\node[font=\small] at (1.6,4.3) {(b) リザバー};
\node[blk,text width=1.1cm] (X) at (-0.4,1.0) {入力\\$u(t)$};
\draw[very thick,fill=orange!10] (1.6,1.0) circle (0.8);
\foreach \a/\r in {20/0.35,80/0.5,150/0.3,210/0.55,280/0.4,330/0.2,110/0.15}{\fill[red!60!black] ({1.6+\r*cos(\a)},{1.0+\r*sin(\a)}) circle (0.05);}
\node[font=\scriptsize] at (1.6,-0.2) {オルガノイド：教師なし};
\node[blk,text width=1.8cm,fill=green!8] (W) at (4.0,1.0) {読み出し\\$W_{\text{out}}$};
\draw[->,thick] (X) -- (0.8,1.0);
\draw[->,thick] (2.4,1.0) -- node[above,font=\scriptsize] {$\mathbf r(t)$} (W);
\draw[->,thick] (W) -- (4.0,2.3) node[above,font=\scriptsize] {$y(t)$};
\node[font=\scriptsize,green!40!black] at (4.0,0.3) {教師あり学習};
\end{scope}
\end{tikzpicture}
\caption{生きたネットワークに計算させる二つの方法。(a) 閉ループ：ボールの位置は電流として与えられ、運動領域のスパイクがラケットを動かし、打った後の予測可能な刺激またはランダムな刺激が教師となる。(b) リザバー\eqref{eq:reservoir}：オルガノイドには学習信号を与えない。オルガノイドは入力を、記憶をもつ多数の非線形応答に分解する。誤差から学習するのは外側の線形読み出し一つだけである。ただしオルガノイド自身の結合もこのとき変化する。教師なしで。}
\label{fig:livingmachine}
\end{figure}

\section{生きたリザバー}

生きたネットワークの使い方はもう一つある。まったく教えないことである。工学ではこれを\emph{リザバー計算}と呼ぶ\cite{Maass2002,JaegerHaas2004}。記憶をもつ出来合いの非線形システムを持ってくる。入力への応答が豊かで、直近の過去に依存しさえすれば何でもよい。入力 $u(t)$ がそれを揺り動かし、応答ベクトル $\mathbf r(t)$ が得られる。学習するのは線形の読み出し
\begin{empheq}[box=\keybox]{equation}
y(t) \;=\; W_{\text{out}}\,\mathbf r(t) ,
\label{eq:reservoir}
\end{empheq}
だけであり、その重みは第\ref{sec:motorlearning}章の最小二乗則\eqref{eq:lms}で決める。リザバーは、第\ref{sec:motorlearning}章の小さな\nt{GLUT}ニューロンがしていたことをする。入力を長いベクトルに分解し、目的のクラスが平面で分けられるようにするのである（第\ref{sec:dendrites}章）。

電極アレイ上のオルガノイドは、こうしたリザバーの役に向いていた。電流パターンとして与えた音声は、線形読み出しを一つ学習させたあと、話者を約 $78\%$ の精度で区別できた\cite{Cai2023}。$78\%$ という精度は著者ら自身が報告し、報道が繰り返した数字である。有用な結果だが、その中のどこに「知恵」があるかを理解しておくことが大切である。オルガノイドは非線形性と減衰する記憶を与え、誤差による学習は外側のシリコンで行われる（図\ref{fig:livingmachine}）。ただしオルガノイドは不変のままではない。著者らによれば、訓練の間にオルガノイド自身の機能的結合が変化し、彼らはこれをオルガノイド内の教師なし学習と呼んでいる\cite{Cai2023}。精度のどれだけがこの変化によるもので、どれだけが読み出しによるものかは切り分けられていない。

\section{テストベンチがマシンについて語ること}

\emph{エネルギー。}推定\eqref{eq:energy}によれば、スパイク1回のコストは約 $10^{-10}$~J である。100万個のニューロンからなる培養が1秒に1スパイクの頻度で発火すると、信号伝達に使う電力は
\begin{empheq}[box=\keybox]{equation}
P \;\approx\; 10^{6}\times1\,\text{s}^{-1}\times10^{-10}\,\text{J} \;=\; 10^{-4}\,\text{W},
\label{eq:dishpower}
\end{empheq}
すなわち0.1ミリワットである。これらのスパイクを刺激し、記録し、処理する電子回路は、桁違いに多くを使う。第\ref{sec:debugging}章と同じく、ボトルネックは計算ではなく、計算との通信である。

\emph{足りない部品。}テストベンチは、制御レジスタなしで、\nt{GLUT}のバスと\nt{GABA}のフロー制御にどれだけのことができるかを示す。自己組織化し、一斉発火を生み、フィードバックによって応答を変え、よいリザバーになる。それらなしにできないのは、何を成功とみなすかを自分で決めることである。この信号を、テストベンチではエンジニアが与え、脳では、第\ref{sec:dofa}--\ref{sec:acet}章で見たように、ブロードキャストチャネルが与える。

\emph{倫理。}オルガノイドはヒトの細胞から育てられ、複雑になるほど、そうした組織が何かを感じうるかという問いは重くなる。工学的な見方は部分的な答えを示す。第\ref{sec:pain}章の痛みは一つのセンサの信号ではなく、ゲート、音量調節器、ブロードキャストチャネルを備えた警報システム全体であり、それらはテストベンチにはない。しかしこれは推論であって証明ではない。この分野自身も、倫理的評価を実験と並行して、最初から行うことを提案している\cite{Smirnova2023}。

\begin{keyblock}
\begin{proposition}[テストベンチ上の生きたネットワーク]
\label{prop:livingmachine}
電極アレイ上の\nt{GLUT}ニューロンと\nt{GABA}ニューロンのネットワークは、自ら一斉発火\eqref{eq:burstinterval}を生み、フィードバックによって応答を変え、外で学習させた読み出しのためのリザバー\eqref{eq:reservoir}になりうる。これらはすべて測定済みである。そうしたネットワークが何らかの一般的な規則、たとえば自分の入力を予測するという規則で学習するというのは仮説であり、その「知性」をめぐる大げさな言葉を大多数の専門家は受け入れていない。
\end{proposition}
\end{keyblock}

\section{光のフラッシュで放出するドーパミン}
\label{sec:dishdopamine}

培養にドーパミンを教師として与えた直接の実験として私たちが知っている唯一のものは、研究者が遠隔で接続するオルガノイドのプラットフォームで行われた\cite{Jordan2024}。オルガノイドを浸す液に、光感受性の「ケージ」に閉じ込めたドーパミンを加えた。ケージに入った分子は受容体に作用しない。ケージドドーパミンの濃度は $1$~mM である。ネットワークに報酬を与えたいときは紫外線を当てる。ケージが壊れ、ドーパミンが空間に出ていく。

ループは次のようにできている。同じ刺激を何度も与え、それが以前より多くのスパイクを引き起こしたら、フラッシュを点けてドーパミンを放出する。$240$ 回の繰り返しで、誘発スパイクの数は全体として増えた。著者ら自身、この実験一つだけでは増加がドーパミンによるとは結論できないと書いている\cite{Jordan2024}。

エンジニアにとって、これはシャーレの中に三要素則\eqref{eq:threefactor}を組む最初の試みである。送り手と受け手の活動が適格度トレースを残し、共通の信号が変化を定着させるかどうかを決める。しかしここでの信号は正にしかならない。報酬があるかないかであり、負の誤差 $\delta<0$ をこの装置は出せない。そのうえドーパミンは広がり、ゆっくりとしか除去されない。\eqref{eq:lowpass}の濃度と同じである。そのため頻繁なフラッシュは単に全体のレベルを上げるだけかもしれない。学習をこれと区別するには二つの対照が必要である。ネットワークの応答と無関係なランダムな時刻に同じ数のフラッシュを与える対照と、ケージドドーパミンなしで同じフラッシュを与える対照である。さらに休止後の確認も必要である。ドーパミンが去ったあとも変化が残っているかどうかである。

\section{ドーパミンがシリコンを教える}
\label{sec:biohybrid}

逆向きの実験では、生きた細胞が電子回路を教える\cite{Keene2020}。ドーパミンを分泌する細胞を、有機ニューロモルフィック素子の上の微小流路で育てる。この素子はトランジスタで、そのチャネルのコンダクタンスがシナプスの重みの役をする。細胞から出たドーパミンは素子の電極で酸化され、それがコンダクタンスを変える。この装置はシナプス間隙から伝達物質を除去する仕組みも再現しているので、重みは長く変えておくことも、元に戻すこともできる。

回路としては、これは化学入力の漏れ積分器である。重みはドーパミンの流入を蓄積し、「汲み出し」が忘れる速さを決める。\eqref{eq:lowpass}と同じ RC 回路である。ここで肝心なのは接続の向きである。第\ref{sec:debugging}章のプローブにブロードキャストレジスタが見えないのは、それが化学的だからである。この素子はまさにその化学レジスタを読み取り、電気的な重みに変える。ブレイン・コンピュータ・インターフェースにとって、これはデータバスだけでなく制御レジスタも聴き取れるセンサーへの道である。

\section{自前のドーパミンニューロンをもつオルガノイド}
\label{sec:midbrainorg}

ヒトの幹細胞から、\nt{DOPA} ニューロンのある脳領域に似たオルガノイドを育てることができる。そこには働くドーパミンニューロンがあり、ドーパミンを分泌する\cite{Jo2016}。こうしたオルガノイドは主に病気の研究のために育てられている。そのドーパミンを別のネットワークの教師として使った実験は、私たちには見つからなかった。

生きたニューロンでできたマシンにとって、これは足りない部品である。第\ref{sec:livingmachine}章のテストベンチは、制御レジスタなしのデータバスとフロー制御だった。ここにはブロードキャスト信号の源がすでにある。足りないのは批評家、つまり予測誤差\eqref{eq:ctd}を計算してこれらのニューロンを制御するネットワークである。大脳皮質のオルガノイドをこうしたオルガノイドとつなぎ、誤差に応じてドーパミンが放出されるようにすることは未解決の課題であり、今のところ実験ではなく仮説である。

\section{ドーパミンなしで棒を立てるオルガノイド}
\label{sec:cartpole}

最近の実験で最も完全なものは、ドーパミンをまったく使わない\cite{Robbins2026}。マウス大脳皮質のオルガノイドを、「台車の上の棒」（倒立振子）の課題と閉ループでつないだ。オルガノイドの活動が台車を動かし、棒の位置は刺激として戻される。試行の合間に、オルガノイドは短い高頻度の学習刺激を受ける。それをどれにするかは強化学習プログラムが選ぶ。

結果は測定されている。
\begin{itemize}
\item 大多数のオルガノイドで、プログラムが選んだ学習信号は、ランダムに選んだ信号や信号なしより良い成績をもたらした。
\item $45$ 分の休止後、改善は残らなかった。
\item AMPA と NMDA の両方のグルタミン酸受容体を遮断すると、改善はなくなった。
\end{itemize}

最後の項目は、学習したものがどこにあるかを語る。\nt{GLUT} シナプスの中であり、\ref{sec:weightstore}節で入力と出力の同時検出器として働く受容体を通じてである。2番目の項目は何が足りないかを示す。速い変化はあるが定着がない。第\ref{sec:motorlearning}章の、遅い学習者を欠いた速い学習者と同じである。そして教師の役割は変わった。電流パターンを選ぶエンジニアはプログラムに置き換わったが、教師は依然としてシャーレの外にいる。

電極アレイ上の培養を学習させたもっと古い実験の中にも、ドーパミンを使ったものは一つも見つからなかった。学習信号はどれも電気刺激だった。

\section{誰が培養を教えるのか}

\begin{table}[t]
\centering
\footnotesize
\setlength{\tabcolsep}{4pt}
\renewcommand{\arraystretch}{1.15}
\begin{tabular}{>{\raggedright}p{2.7cm}>{\raggedright}p{2.5cm}>{\raggedright}p{3.2cm}>{\raggedright\arraybackslash}p{4.4cm}}
\toprule
実験 & 教えるのは誰か & 学習信号 & わかっていること \\
\midrule
目的の応答で刺激を止める & エンジニア & 刺激の終了 & 応答が定着する（測定済み） \\
DishBrain（第\ref{sec:dishbrain}章） & エンジニア & 予測可能な電流またはランダム電流 & セッション内のラリーの有意な伸びは完全なフィードバックのときだけ、無音では効果が小さい（測定済み）。セッション間では不安定。原因は議論中 \\
台車の上の棒 & 強化学習プログラム & プログラムが選ぶ高頻度刺激 & ランダムより良いが、休止後に残らない（測定済み） \\
フラッシュで放出するドーパミン & 「スパイクが増えたら報酬」という規則 & 光で解放されるドーパミン & スパイクの増加は観察。ドーパミンの役割は未証明 \\
バイオハイブリッドシナプス & 生きた細胞 & 細胞からのドーパミン & 素子の重みの長期変化と復帰（測定済み） \\
\nt{DOPA} ニューロンをもつオルガノイド & --- & --- & ドーパミン源はある。教師としては未使用 \\
\bottomrule
\end{tabular}
\caption{チップ上の生きたニューロンを誰が何で教えるか。}
\label{tab:dishteachers}
\end{table}

培養自体が学習するどの実験でも、教師は今のところ外にいる（表\ref{tab:dishteachers}）。エンジニアか、プログラムか、エンジニアが定めた規則である。ドーパミンがシャーレに入ったのは一度だけで、その役割はまだ証明されていない。第\ref{sec:dofa}章のように批評家、誤差、適格度トレースを備えた本物のブロードキャストチャネルをシャーレの中に組むことが次の一歩であり、まだ誰も踏み出していない。


\section{まとめ}

\begin{table}[t]
\centering
\footnotesize
\setlength{\tabcolsep}{4pt}
\renewcommand{\arraystretch}{1.15}
\begin{tabular}{>{\raggedright}p{2.8cm}>{\raggedright}p{4.2cm}>{\raggedright\arraybackslash}p{5.8cm}}
\toprule
実験 & 生きたネットワークがすること & わかっていること \\
\midrule
入力のないネットワーク & 間隔1秒から数分の一斉発火\eqref{eq:burstinterval} & 測定済み。シナプスの枯渇による機構は広く受け入れられたモデル \\
黙る教師 & 刺激を止めると目的の応答が定着する & 測定済み \\
テニスゲーム & ラリーが長くなる。セッション内の有意な伸びは完全なフィードバックのときだけ & 振る舞いの変化は測定済み。セッション間の学習は不安定。効果の大きさと入力予測による説明には異論がある \\
仮想の体 & 刺激パターンを選べば目標へ動く & 測定済み \\
リザバー & 読み出し\eqref{eq:reservoir}のための非線形性と記憶 & 測定済み。誤差による学習は外側のシリコンで。オルガノイドの結合も変化する \\
エネルギー & 100万ニューロンあたり約 $10^{-4}$~W~\eqref{eq:dishpower} & \eqref{eq:energy}による推定 \\
\bottomrule
\end{tabular}
\caption{生きたニューロンでできたマシンが示したこと。}
\label{tab:livingmachine}
\end{table}

生きたニューロンでできたマシン（表\ref{tab:livingmachine}）は、制御レジスタなしのデータバスとフロー制御に何ができるかを見せるテストベンチである。できることは多いが、何が良いかを自分で決めることはできない。テストベンチではエンジニアが自分の電流パターンでこの役を担い、脳では、本書の中盤で扱ったわずかなブロードキャスト配線が担う。

\section{問題}

\begin{exercise}[\lvlA]
\label{ex:21:1}
\emph{一斉発火の間隔。}一斉発火は蓄えをゼロまで使い果たし、その後 $\tau_{\text{rec}}\,\dd x/\dd t=1-x$ に従って回復し、$x=x^*$ で次の一斉発火が始まる。$\tau_{\text{rec}}=0.8$~s。(a)~$x^*=0.9$ と $0.99$ のときの間隔を求めよ。(b)~閾値 $x^*$ が $0.99$ の周りで $\pm0.005$ ランダムに変わる。間隔はどの範囲に入るか。
\end{exercise}

\begin{exercise}[\lvlA]
\label{ex:21:2}
\emph{テストベンチのエネルギー。}(a)~推定\eqref{eq:dishpower}は脳全体（$8.6\cdot10^{10}$ ニューロン）でどれだけの電力を与えるか。(b)~テストベンチは $1024$ 個の電極を $20$~kHz、1サンプル $10$~bit で記録する。1 bit あたり $1$~nJ のとき、このデータの伝送は培養自体の電力の何倍か。
\end{exercise}

\begin{exercise}[\lvlB]
\label{ex:21:3}
\emph{設計：一斉発火を抑える。}多数の電極からのまばらな刺激で、各シナプスが頻度 $r_b$ で伝達物質を放出する。蓄えは $\tau_{\text{rec}}\,\dd x/\dd t=1-x-Ur_b\tau_{\text{rec}}x$、$U=0.5$、$\tau_{\text{rec}}=0.8$~s。蓄えが $x^*=0.9$、$0.99$ に届かなくなる最小の $r_b$ はいくらか。そのときシナプスはどれだけ弱まるか。
\end{exercise}

\begin{exercise}[\lvlB]
\label{ex:21:4}
\emph{リザバーの記憶はニューロン数を超えない。}出力 $\mathbf r(t)$ が $N$ 個のリザバーに、平均ゼロ・分散1の白色入力 $u(t)$ を与える。$m_k$ を読み出し $\mathbf w_k\cdot\mathbf r(t)$ と $u(t-k)$ の相関の2乗の最大値とする。記憶容量 $\mathrm{MC}=\sum_{k\ge0}m_k\le N$ を証明せよ。
\end{exercise}

\begin{exercise}[\lvlB]
\label{ex:21:5}
\emph{シミュレーション：シリコンのリザバー。}リザバー $\mathbf r(t+1)=\tanh(W\mathbf r(t)+\mathbf v\,u(t)+\mathbf c)$、$N=30$、$W$ はスペクトル半径 $0.9$ のランダム行列、$v_i\sim U(-0.5,0.5)$、$c_i\sim U(-0.2,0.2)$、$u\sim U(-1,1)$、読み出しはリッジ回帰。問題\ref{ex:21:4}の記憶容量を $\tanh$ ありとなしで求めよ。どちらのリザバーがより多く覚えているか。
\end{exercise}

\begin{exercise}[\lvlB]
\label{ex:21:6}
\emph{生きたリザバーの読み出し。}オルガノイドは $1024$ チャネルと、2クラスの訓練記録 $P=240$ 個を与える。(a)~問題\ref{ex:19:3}の式で、ランダムなラベル付けが分離可能になる確率を $1\%$ 以下に抑えるには、特徴量 $n$ をいくつまで残せるか。(b)~$100$ 個の検証記録で精度 $78\%$ は、当て推量より何シグマ高いか。
\end{exercise}

\begin{exercise}[\lvlC]
\label{ex:21:7}
\emph{研究：ブロードキャストチャネルのない教師。}第\ref{sec:dofa}章の三要素則を $8$--$1024$ 個の電極でテストベンチ上に実現する刺激プロトコルと、読み出しを凍結して重みの学習とネットワーク状態のずれを区別する対照実験を提案せよ。どんな結果が学習の仮説を反証するか。
\end{exercise}
